\chapter{Fully Nonlinear Equations}

Let $\Omega\subset\mathbb R^n$ and let $\mathbb R^{n\times n}$ denote the
space of real symmetric $n\times n$ matrices.  A general second-order equation
has the form
\begin{equation}\tag{17.1}
F[u]=F(x,u,Du,D^2u)=0,
\end{equation}
where $F$ is defined on
$\Gamma=\Omega\times\mathbb R\times\mathbb R^n\times\mathbb R^{n\times n}$.
Write
\[
F_{ij}(x,z,p,r)=\frac{\partial F}{\partial r_{ij}}(x,z,p,r).
\]
On a subset $\mathcal U\subset\Gamma$, the operator is \emph{elliptic} when
$[F_{ij}]$ is positive definite, \emph{uniformly elliptic} when the ratio of
its largest and smallest eigenvalues $\Lambda/\lambda$ is bounded, and
\emph{strictly elliptic} when $1/\lambda$ is bounded.  These terms are applied
\emph{with respect to} $u\in C^2(\Omega)$ when the corresponding condition
holds on the range of $x\mapsto(x,u,Du,D^2u)$.

The Monge--Amp\`ere equation and the prescribed Gauss-curvature equation are
\begin{equation}
\det D^2u-f(x)=0.
\tag{17.2}
\end{equation}
\begin{equation}
\det D^2u-K(x)(1+|Du|^2)^{(n+2)/2}=0.
\tag{17.3}
\end{equation}
More generally, an equation of Monge--Amp\`ere type is
\begin{equation}\tag{17.4}
\det D^2u-f(x,u,Du)=0,
\end{equation}
with $f>0$.  These equations are elliptic on the cone of positive Hessians,
so their elliptic solutions are convex; for $F(r)=\det r$, the matrix
$[F_{ij}(r)]$ is the cofactor matrix of $r$.

For $0<\alpha\le 1/n$, let $\mathcal L_\alpha$ consist of the operators
$Lu=a^{ij}D_{ij}u$ whose bounded measurable coefficients satisfy
\[
a^{ij}\xi_i\xi_j\ge\alpha|\xi|^2,
\qquad \sum_i a^{ii}=1.
\]
If $\mathcal C_1(r)$ and $\mathcal C_n(r)$ are the least and greatest
eigenvalues of $r$, the Pucci extremal operators are
\begin{equation}
\tag{17.5}
\begin{aligned}
M_\alpha[u]&=\sup_{L\in\mathcal L_\alpha}Lu
&=\alpha\Delta u+(1-n\alpha)\mathcal C_n(D^2u),\\
m_\alpha[u]&=\inf_{L\in\mathcal L_\alpha}Lu
&=\alpha\Delta u+(1-n\alpha)\mathcal C_1(D^2u).
\end{aligned}
\end{equation}
and $M_\alpha[-u]=-m_\alpha[u]$.  The equations
\begin{equation}\tag{17.6}
M_\alpha[u]=f,
\qquad m_\alpha[u]=f
\end{equation}
are uniformly elliptic in the Lipschitz sense described below.

For a family indexed by $\nu\in V$, set
\begin{equation}\tag{17.7}
L_\nu u=a_\nu^{ij}(x)D_{ij}u+b_\nu^i(x)D_i u+c_\nu(x)u.
\end{equation}
The associated Bellman equation is
\begin{equation}\tag{17.8}
F[u]=\inf_{\nu\in V}(L_\nu u-f_\nu)=0.
\end{equation}
For an operator locally Lipschitz in $r$, ellipticity means that $[F_{ij}]$ is
positive wherever it exists.  In particular, (17.8) is elliptic if positive
functions $\lambda,\Lambda$ satisfy
\begin{equation}\tag{17.9}
\lambda(x)|\xi|^2\le a_\nu^{ij}(x)\xi_i\xi_j
\le\Lambda(x)|\xi|^2
\end{equation}
for every $x\in\Omega$, $\nu\in V$, and $\xi\in\mathbb R^n$; it is uniformly
elliptic when $\Lambda/\lambda\in L^\infty(\Omega)$.

\section*{17.1. Maximum and Comparison Principles}

\begin{theorem}[Comparison principle for fully nonlinear equations]\label{gt-ch17-thm-17-1}\dcref{ch17:17.1}\group{chapter-17}
\source{gilbarg-trudinger:page-0455}
Let $u,v\in C^0(\overline\Omega)\cap C^2(\Omega)$ satisfy
$F[u]\ge F[v]$ in $\Omega$ and $u\le v$ on $\partial\Omega$.  Assume that
\begin{enumerate}
\item[(i)] $F$ is continuously differentiable in $z,p,r$ on $\Gamma$;
\item[(ii)] $F$ is elliptic with respect to
$\theta u+(1-\theta)v$ for every $\theta\in[0,1]$;
\item[(iii)] $F$ is nonincreasing in $z$ for fixed $(x,p,r)$.
\end{enumerate}
Then $u\le v$ in $\Omega$.
\end{theorem}
\begin{proof}
\sketch Put $w=u-v$ and $u_\theta=\theta u+(1-\theta)v$.  The identity
\[
F[u]-F[v]=a^{ij}D_{ij}w+b^iD_iw+cw,
\]
where $a^{ij}$, $b^i$, and $c$ are the integrals from $0$ to $1$ of
$F_{ij}$, $F_{p_i}$, and $F_z$ along $u_\theta$, reduces the claim to the
weak maximum principle: $[a^{ij}]$ is positive and $c\le0$.
\end{proof}

When the strong maximum principle applies, the conclusion sharpens to the
alternative that $u<v$ throughout $\Omega$ or $u=v$ throughout $\Omega$.

\begin{corollary}[Dirichlet uniqueness]\label{gt-ch17-cor-17-2}\dcref{ch17:17.2}\group{chapter-17}
\source{gilbarg-trudinger:page-0456}
Under hypotheses (i)--(iii) of Theorem~\ref{gt-ch17-thm-17-1}, if
$u,v\in C^0(\overline\Omega)\cap C^2(\Omega)$ satisfy $F[u]=F[v]$ in
$\Omega$ and $u=v$ on $\partial\Omega$, then $u=v$ in $\Omega$.
\end{corollary}
\begin{proof}
\sketch Apply Theorem~\ref{gt-ch17-thm-17-1} to $(u,v)$ and to $(v,u)$.
\end{proof}

For $u\in C^2(\Omega)$, the identity
\begin{equation}\tag{17.10}
F[u]=a^{ij}(x,u,Du)D_{ij}u+b(x,u,Du),
\end{equation}
holds with
\[
a^{ij}(x,z,p)=\int_0^1F_{ij}(x,z,p,\theta D^2u(x))\,d\theta,
\qquad b(x,z,p)=F(x,z,p,0).
\]
Set
\[
\mathcal E=F_{ij}p_ip_j,\quad \mathcal E^*=\mathcal E/|p|^2,
\qquad \mathcal D=\det[F_{ij}],\quad \mathcal D^*=\mathcal D^{1/n}.
\]

\begin{theorem}[Maximum principle under natural growth]\label{gt-ch17-thm-17-3}\dcref{ch17:17.3}\group{chapter-17}
\source{gilbarg-trudinger:page-0457}
Suppose $F$ is elliptic in $\Omega$ and there are constants
$\mu_1,\mu_2\ge0$ such that, with either denominator indicated,
\begin{equation}\tag{17.11}
\frac{F(x,z,p,0)\operatorname{sign}z}{\mathcal E^*\ \text{(or }\mathcal D^*\text{)}}
\le\mu_1|p|+\mu_2
\qquad ((x,z,p,r)\in\Gamma).
\end{equation}
If $u\in C^0(\overline\Omega)\cap C^2(\Omega)$ satisfies $F[u]\ge0$, then
\begin{equation}\tag{17.12}
\sup_\Omega u\le\sup_{\partial\Omega}u^++C\mu_2.
\end{equation}
If $F[u]=0$, the same formula holds with $u$ and $u^+$ replaced by $|u|$.
Here $C=C(\mu_1,\operatorname{diam}\Omega)$.
\end{theorem}

\begin{theorem}[Integral maximum bound for Monge--Amp\`ere equations]\label{gt-ch17-thm-17-4}\dcref{ch17:17.4}\group{chapter-17}
\source{gilbarg-trudinger:page-0457}
Let $F$ be the Monge--Amp\`ere operator (17.4).  Suppose there are
$g\in L^1_{\mathrm{loc}}(\mathbb R^n)$, $h\in L^1(\Omega)$, both
nonnegative, and $N\ge0$ such that
\begin{equation}
|f(x,z,p)|\le \frac{h(x)}{g(p)}
\quad (x\in\Omega,\ |z|\ge N,\ p\in\mathbb R^n).
\tag{17.13}
\end{equation}
\begin{equation}
\int_\Omega h\,dx<\int_{\mathbb R^n}g(p)\,dp=:g_\infty.
\tag{17.14}
\end{equation}
If $u\in C^0(\overline\Omega)\cap C^2(\Omega)$ satisfies $F[u]\ge0$, then
\begin{equation}\tag{17.15}
\sup_\Omega u\le\sup_{\partial\Omega}u^+
+C\operatorname{diam}\Omega+N,
\end{equation}
and for $F[u]=0$ the same estimate holds with absolute values.  The constant
$C$ depends only on $g$ and $h$.  For (17.2),
\begin{equation}\tag{17.16}
\sup_\Omega u\le\sup_{\partial\Omega}u^+
+\frac{\operatorname{diam}\Omega}{\omega_n^{1/n}}
\left(\int_\Omega|f|\right)^{1/n},
\end{equation}
again with absolute values in the equality case.  For (17.3), (17.15) holds
with $C=C(n,K_0)$ provided
\begin{equation}\tag{17.17}
K_0=\int_\Omega|K(x)|\,dx<\omega_n.
\end{equation}
\end{theorem}

\begin{theorem}[Lipschitz comparison principle]\label{gt-ch17-thm-17-5}\dcref{ch17:17.5}\group{chapter-17}
\source{gilbarg-trudinger:page-0458}
Let $u,v\in C^0(\overline\Omega)\cap C^2(\Omega)$ satisfy
$F[u]\ge F[v]$ in $\Omega$ and $u\le v$ on $\partial\Omega$.  Assume that
\begin{enumerate}
\item[(i)] $F$ is locally uniformly Lipschitz in $z,p,r$ on $\Gamma$;
\item[(ii)] $F$ is elliptic in $\Omega$;
\item[(iii)] $F$ is nonincreasing in $z$ for fixed $(x,p,r)$;
\item[(iv)] $|F_p|/\lambda$ is locally bounded on $\Gamma$.
\end{enumerate}
Then $u\le v$ in $\Omega$.
\end{theorem}
\begin{proof}
\sketch Linearize $F[u]-F[v]$ along the segment between $u$ and $v$ at
points of differentiability.  The Lipschitz hypothesis supplies almost-everywhere
coefficients, ellipticity gives the principal part, and (iii)--(iv) permit the
maximum principle for the resulting linear inequality.
\end{proof}

For the Bellman operator (17.8), Theorem~\ref{gt-ch17-thm-17-5} applies when (17.9) holds,
$c_\nu\le0$, and $|b_\nu|/\lambda$ is locally bounded uniformly in $\nu$.
The conclusions of this section also hold for strong solutions in
$W^{2,n}_{\mathrm{loc}}(\Omega)$.

\section*{17.2. The Method of Continuity}

For Banach spaces $\mathcal B_1,\mathcal B_2$, a map $F$ is Fr\'echet
differentiable at $u$ with derivative $F_u$ when
\begin{equation}\tag{17.18}
\frac{\|F[u+h]-F[u]-F_uh\|_{\mathcal B_2}}{\|h\|_{\mathcal B_1}}
\longrightarrow0.
\end{equation}
It is continuously differentiable when it is differentiable on a neighbourhood
and $v\mapsto F_v$ is continuous as a map into the bounded linear operators.
The chain and mean-value estimates are
\begin{equation}
(G\circ F)_u=G_{F[u]}\circ F_u.
\tag{17.19}
\end{equation}
\begin{equation}
\|F[u]-F[v]\|_{\mathcal B_2}
\le \sup_{w\in[u,v]}\|F_w\|\,\|u-v\|_{\mathcal B_1}.
\tag{17.20}
\end{equation}

\begin{theorem}[Banach-space implicit function theorem]\label{gt-ch17-thm-17-6}\dcref{ch17:17.6}\group{chapter-17}
\source{gilbarg-trudinger:page-0460}
Let $\mathcal B_1,\mathcal B_2,X$ be Banach spaces and let $G$ map an open
subset of $\mathcal B_1\times X$ into $\mathcal B_2$.  Suppose
\begin{enumerate}
\item[(i)] $G[u_0,\sigma_0]=0$;
\item[(ii)] $G$ is continuously differentiable at $(u_0,\sigma_0)$;
\item[(iii)] the partial derivative $L=G^1_{(u_0,\sigma_0)}$ is invertible.
\end{enumerate}
Then a neighbourhood $\mathcal N$ of $\sigma_0$ exists such that for each
$\sigma\in\mathcal N$ there is $u_\sigma\in\mathcal B_1$ with
$G[u_\sigma,\sigma]=0$.
\end{theorem}
\begin{proof}
\sketch Rewrite the equation as
$u=T_\sigma u:=u-L^{-1}G[u,\sigma]$.  Equations (17.19)--(17.20) make
$T_\sigma$ a contraction on a sufficiently small closed ball about $u_0$ for
$\sigma$ close to $\sigma_0$.  The fixed-point map is differentiable at
$\sigma_0$, with derivative $-L^{-1}G^2_{(u_0,\sigma_0)}$.
\end{proof}

Let $F:\mathcal U\subset\mathcal B_1\to\mathcal B_2$, fix
$\psi\in\mathcal U$, and define
\begin{equation}\tag{17.21}
\begin{aligned}
G[u,t]&=F[u]-tF[\psi],\\
S&=\{t\in[0,1]:G[u,t]=0\text{ for some }u\in\mathcal U\},\\
E&=\{u\in\mathcal U:G[u,t]=0\text{ for some }t\in[0,1]\}.
\end{aligned}
\end{equation}

\begin{theorem}[Method of continuity criterion]\label{gt-ch17-thm-17-7}\dcref{ch17:17.7}\group{chapter-17}
\source{gilbarg-trudinger:page-0460}
Assume that $F$ is continuously differentiable on $E$, that each derivative
$F_u$ for $u\in E$ is invertible, and that $S$ is closed in $[0,1]$.  Then
$F[u]=0$ is solvable for some $u\in\mathcal U$.
\end{theorem}
\begin{proof}
\sketch The set $S$ is nonempty because $1\in S$.  Theorem~\ref{gt-ch17-thm-17-6} makes it
open, while it is closed by hypothesis; hence $S=[0,1]$, and $0\in S$ gives
the solution.
\end{proof}

For (17.1) acting from $C^{2,\alpha}(\overline\Omega)$ to
$C^\alpha(\overline\Omega)$, its derivative is
\begin{equation}\tag{17.22}
F_uh=F_{ij}D_{ij}h+F_{p_i}D_i h+F_zh,
\end{equation}
with every coefficient evaluated at $(x,u,Du,D^2u)$.

\begin{theorem}[Dirichlet method of continuity]
\label{gt-ch17-thm-dirichlet-continuity}
\dcref{ch17:17.8}\group{chapter-17}\source{gilbarg-trudinger:page-0461}
Let $\Omega$ be bounded with $\partial\Omega\in C^{2,\alpha}$,
$0<\alpha<1$, let $\mathcal U\subset C^{2,\alpha}(\overline\Omega)$ be open,
and let $\phi\in\mathcal U$.  Set
\[
E=\{u\in\mathcal U:F[u]=\sigma F[\phi]\text{ for some }\sigma\in[0,1],
\ u=\phi\text{ on }\partial\Omega\}.
\]
Suppose $F\in C^{2,\alpha}(\overline\Gamma)$ and
\begin{enumerate}
\item[(i)] $F$ is strictly elliptic with respect to every $u\in E$;
\item[(ii)] $F_z(x,u,Du,D^2u)\le0$ for every $u\in E$;
\item[(iii)] $E$ is bounded in $C^{2,\alpha}(\overline\Omega)$;
\item[(iv)] $\overline E\subset\mathcal U$.
\end{enumerate}
Then $F[u]=0$ in $\Omega$, $u=\phi$ on $\partial\Omega$, has a solution in
$\mathcal U$.
\end{theorem}
\begin{proof}
\sketch Replace $u$ by $u-\phi$ and apply Theorem~\ref{gt-ch17-thm-17-7} to
$G[u,\sigma]=F[u+\phi]-\sigma F[\phi]$ on the zero-boundary subspace.
Strict ellipticity and $F_z\le0$ make (17.22) invertible.  The bound on $E$,
the compact embedding, and $\overline E\subset\mathcal U$ make the
continuation set closed.
\end{proof}

\section*{17.3. Equations in Two Variables}

Differentiating (17.1) gives, for $k=1,\ldots,n$,
\begin{equation}\tag{17.23}
F_{ij}D_{ijk}u+F_{p_i}D_{ik}u+F_zD_ku+F_{x_k}=0.
\end{equation}
In dimension two assume, along $u$, that
\begin{equation}\tag{17.24}
0<\lambda|\xi|^2\le F_{ij}\xi_i\xi_j\le\Lambda|\xi|^2,
\qquad |F_p|,|F_z|,|F_x|\le\lambda\mu.
\end{equation}

\begin{theorem}[Interior Hessian H\"older estimate in dimension two]
\label{gt-ch17-thm-interior-hessian-holder-two-dimensional}
\dcref{ch17:17.9}\group{chapter-17}\source{gilbarg-trudinger:page-0463}
Let $u\in C^3(\Omega)$ solve $F[u]=0$ in $\Omega\subset\mathbb R^2$.
Assume $F\in C^1(\Gamma)$ is elliptic with respect to $u$ and satisfies
(17.24).  For every $\Omega'\Subset\Omega$,
\begin{equation}\tag{17.25}
[D^2u]_{\alpha;\Omega'}\le\frac{C}{d^\alpha}
\bigl(|D^2u|_{0;\Omega}+\mu d(1+|Du|_{1;\Omega})\bigr),
\end{equation}
where $d=\operatorname{dist}(\Omega',\partial\Omega)$ and $C,\alpha$ depend
only on $\Lambda/\lambda$.
\end{theorem}
\begin{proof}
\sketch Each $D_ku$ satisfies the linear uniformly elliptic equation (17.23).
The two-dimensional H\"older gradient estimate for that equation, applied on
balls of radius comparable to $d$, gives (17.25).
\end{proof}

\begin{theorem}[Global Hessian H\"older estimate in dimension two]
\label{gt-ch17-thm-global-hessian-holder-two-dimensional}
\dcref{ch17:17.10}\group{chapter-17}\source{gilbarg-trudinger:page-0464}
Let $u\in C^3(\Omega)\cap C^2(\overline\Omega)$ solve $F[u]=0$ in a domain
$\Omega\subset\mathbb R^2$, with $u=\phi$ on $\partial\Omega$.  Assume
$F\in C^1(\Gamma)$ is elliptic with respect to $u$, (17.24) holds,
$\partial\Omega\in C^3$, and $\phi\in C^3(\overline\Omega)$.  Then
\begin{equation}\tag{17.28}
[D^2u]_{\alpha;\Omega}\le
C\bigl((1+\mu)(1+|u|_{2;\Omega})+|\phi|_{3;\Omega}\bigr),
\end{equation}
where $C,\alpha$ depend only on $\Lambda/\lambda$ and $\Omega$.
\end{theorem}
\begin{proof}
\sketch Flatten the boundary by a $C^3$ coordinate map.  Tangential
derivatives of $u$ satisfy transformed versions of (17.23), with boundary
data differentiated from $\phi$. If
$w=(\partial x_l/\partial y_k)D_lu$, then
\begin{equation}
\tag{17.26}
\begin{aligned}
F_{ij}D_{ij}w+F_{p_i}D_iw+F_zw
&=2F_{ij}D_i\left(\frac{\partial x_l}{\partial y_k}\right)D_{jl}u\\
&\quad+F_{ij}D_{ij}\left(\frac{\partial x_l}{\partial y_k}\right)D_lu
+F_{p_i}D_i\left(\frac{\partial x_l}{\partial y_k}\right)D_lu
-F_{x_l}\frac{\partial x_l}{\partial y_k}.
\end{aligned}
\end{equation}
Boundary gradient estimates and the
equation for the normal derivative give
\begin{equation}
\tag{17.27}
\int_{B_R\cap D^+}|D_y^3u|^2\,dy
\leq CR^{n-2+2\alpha}
\bigl((1+\mu)(1+|Du|_{1;\Omega})+|D^3\phi|_{0;\Omega}^2\bigr).
\end{equation}
Morrey's estimate and
this boundary estimate with
Theorem~\ref{gt-ch17-thm-interior-hessian-holder-two-dimensional}.
\end{proof}

Consequently, in hypothesis (iii) of
Theorem~\ref{gt-ch17-thm-dirichlet-continuity}, a uniform
$C^2(\overline\Omega)$ bound suffices in dimension two when
$\partial\Omega\in C^3$ and $\phi\in C^3(\overline\Omega)$.

The natural two-dimensional structure conditions used below are
\begin{equation}\tag{17.29}
\begin{aligned}
0<\lambda|\xi|^2&\le F_{ij}(x,z,p,r)\xi_i\xi_j\le\Lambda|\xi|^2,\\
|F_p|,|F_z|&\le\mu\lambda,\\
|F_x|&\le\mu\lambda(1+|p|+|r|),
\end{aligned}
\end{equation}
for nonzero $\xi\in\mathbb R^2$ and $(x,z,p,r)\in\Gamma$, where $\lambda$
is nonincreasing and $\Lambda,\mu$ are nondecreasing functions of $|z|$.

\begin{theorem}[A priori $C^{2,\alpha}$ estimates in dimension two]
\label{gt-ch17-thm-c2alpha-estimates-two-dimensional}
\dcref{ch17:17.11}\group{chapter-17}\source{gilbarg-trudinger:page-0464}
Let $u\in C^3(\Omega)$ solve $F[u]=0$ in $\Omega\subset\mathbb R^2$ and
assume (17.29).  Then
\begin{equation}\tag{17.30}
|u|^*_{2,\alpha;\Omega}\le C,
\end{equation}
where $\alpha>0$ depends only on $\Lambda/\lambda$, while $C$ depends on
$\Lambda/\lambda$, $\mu$, and $|u|_{0;\Omega}$.
If additionally $u\in C^3(\overline\Omega)$, $\partial\Omega\in C^3$, and
$u=\phi$ on $\partial\Omega$ for $\phi\in C^3(\overline\Omega)$, then
\begin{equation}\tag{17.31}
|u|_{2,\alpha;\Omega}\le C,
\end{equation}
where $\alpha$ depends on $\Lambda/\lambda$ and $\Omega$, and $C$ depends
also on $\mu$, $\Omega$, and $|u|_{0;\Omega}$.
\end{theorem}
\begin{proof}
\sketch Evaluate the functions $\lambda,\Lambda,\mu$ at
$M=|u|_{0;\Omega}$.  Apply
Theorems~\ref{gt-ch17-thm-interior-hessian-holder-two-dimensional} and
\ref{gt-ch17-thm-global-hessian-holder-two-dimensional}, then absorb the lower
derivative norms by the interior and boundary interpolation inequalities.
\end{proof}

The smoothness hypotheses in Theorems~\ref{gt-ch17-thm-dirichlet-continuity},
\ref{gt-ch17-thm-interior-hessian-holder-two-dimensional}, and
\ref{gt-ch17-thm-global-hessian-holder-two-dimensional} remain valid for
$F\in C^{0,1}(\Gamma)$, $u\in W^{3,2}(\Omega)$, and
$\phi\in W^{3,2}(\Omega)\cap C^{2,\beta}(\overline\Omega)$ when the
corresponding structure conditions hold almost everywhere.

\begin{theorem}[Classical Dirichlet solvability in dimension two]
\label{gt-ch17-thm-dirichlet-solvability-two-dimensional}
\dcref{ch17:17.12}\group{chapter-17}\source{gilbarg-trudinger:page-0465}
Let $\Omega\subset\mathbb R^2$ be bounded with $\partial\Omega\in C^3$ and
let $\phi\in C^3(\overline\Omega)$.  If $F_z\le0$ on $\Gamma$ and $F$
satisfies (17.29), then
\[
F[u]=0\quad\text{in }\Omega,
\qquad u=\phi\quad\text{on }\partial\Omega
\]
has a unique solution $u\in C^{2,\alpha}(\overline\Omega)$ for every
$\alpha<1$.
\end{theorem}
\begin{proof}
\sketch Theorems~\ref{gt-ch17-thm-17-3} and
\ref{gt-ch17-thm-c2alpha-estimates-two-dimensional} supply the uniform
$C^{2,\alpha}$ bound
needed in Theorem~\ref{gt-ch17-thm-dirichlet-continuity}; comparison gives
uniqueness.  Approximation removes
the extra smoothness required for a direct application of the continuity
theorem and regularity upgrades the resulting solution.
\end{proof}

\section*{17.4. H\"older Estimates for Second Derivatives}

For the model equation
\begin{equation}\tag{17.32}
F(D^2u)=g,
\end{equation}
assume $F$ is concave on the range of $D^2u$ and
\begin{equation}\tag{17.33}
\lambda|\xi|^2\le F_{ij}(D^2u)\xi_i\xi_j\le\Lambda|\xi|^2.
\end{equation}
Differentiating in a unit direction $\gamma$ gives
\begin{equation}
\tag{17.34}
F_{ij}D_{ij\gamma}u=D_\gamma g.
\end{equation}
Differentiating again shows that
$w=D_{\gamma\gamma}u$ satisfies
\begin{equation}\tag{17.35}
F_{ij}D_{ij}w\ge D_{\gamma\gamma}g.
\end{equation}

\begin{lemma}[Finite directional decomposition of elliptic matrices]
\label{gt-ch17-lem-finite-directional-decomposition}
\dcref{ch17:17.13}\group{chapter-17}\source{gilbarg-trudinger:page-0466}
Let $S[\lambda,\Lambda]$ be the positive symmetric matrices whose eigenvalues
lie in $[\lambda,\Lambda]$, where $0<\lambda<\Lambda$.  There are unit
vectors $\gamma_1,\ldots,\gamma_N\in\mathbb R^n$ and constants
$0<\lambda^*<\Lambda^*$, depending only on $n,\lambda,\Lambda$, such that
every $\mathcal A=[a^{ij}]\in S[\lambda,\Lambda]$ admits
\begin{equation}\tag{17.38}
\mathcal A=\sum_{k=1}^N\beta_k\gamma_k\otimes\gamma_k,
\qquad \lambda^*\le\beta_k\le\Lambda^*.
\end{equation}
The set of directions may be chosen to contain every $e_i$ and every
$2^{-1/2}(e_i\pm e_j)$ for $i<j$.
\end{lemma}
\begin{proof}
\sketch Every positive matrix belongs to an open simplicial cone generated by
linearly independent dyads; in particular,
\begin{equation}
\tag{17.55}
\mathcal A=\sum_{k=1}^{n(n+1)/2}\gamma_k\otimes\gamma_k.
\end{equation}
These cones cover the
compact set $S[\lambda,\Lambda]$, so a finite subcover gives fixed directions.
Apply the construction to
$\mathcal A-\lambda^*\sum_k\gamma_k\otimes\gamma_k$ for sufficiently small
$\lambda^*>0$ to obtain the lower coefficient bound; compactness gives the
upper bound.  Prescribed finitely many directions can be included at the
start.
\end{proof}

For $w=D_{\gamma\gamma}u$, write
$M_s=\sup_{B_{sR}}w$ and $m_s=\inf_{B_{sR}}w$. The weak Harnack inequality
applied to $M_2-w$ gives
\begin{equation}
\tag{17.36}
\left(R^{-n}\int_{B_R}(M_2-w)^p\right)^{1/p}
\leq C(M_2-M_1+R^2|D^2g|_{0;\Omega}).
\end{equation}
Concavity also gives, for $x,y\in\Omega$,
\begin{equation}
\tag{17.37}
F_{ij}(D^2u(y))(D_{ij}u(x)-D_{ij}u(y))
\geq g(x)-g(y).
\end{equation}
Using (17.38), put $w_l=D_{\gamma_l\gamma_l}u$ and let
$F_{ij}(D^2u(y))=\sum_k\beta_k(y)\gamma_k\otimes\gamma_k$. Then
\begin{equation}
\tag{17.39}
\sum_{k=1}^N\beta_k(y)(w_k(y)-w_k(x))
\leq g(y)-g(x).
\end{equation}
Let
$\omega(R)=\sum_l(\sup_{B_R}w_l-\inf_{B_R}w_l)$. Then
\begin{equation}
\tag{17.40}
\left(R^{-n}\int_{B_R}(w_l-m_{2l})^p\right)^{1/p}
\leq C\bigl(\omega(2R)-\omega(R)
+R|Dg|_{0;\Omega}+R^2|D^2g|_{0;\Omega}\bigr).
\end{equation}

For the model equation (17.32), the same oscillation argument gives, whenever
$B_R\subset B_{R_0}\Subset\Omega$,
\begin{equation}\tag{17.41}
\operatorname{osc}_{B_R}D^2u\le
C\left(\frac R{R_0}\right)^\alpha
\left(\operatorname{osc}_{B_{R_0}}D^2u
+R_0|Dg|_{0;\Omega}+R_0^2|D^2g|_{0;\Omega}\right),
\end{equation}
and hence
\begin{equation}\tag{17.42}
|u|^*_{2,\alpha;\Omega}\le
C\left(|u|^*_{2;\Omega}+|g|^*_{2;\Omega}\right),
\end{equation}
where $C,\alpha$ depend only on $n,\lambda,\Lambda$.

For the general equation (17.1), the corresponding assumptions are
\begin{equation}\tag{17.43}
\lambda|\xi|^2\le F_{ij}(x,u,Du,D^2u)\xi_i\xi_j\le\Lambda|\xi|^2
\end{equation}
and concavity of $F$ in $r$ on the range of $D^2u$.

\begin{theorem}[Interior Hessian H\"older estimate for concave equations]
\label{gt-ch17-thm-interior-evans-krylov}
\dcref{ch17:17.14}\group{chapter-17}\source{gilbarg-trudinger:page-0473}
Let $u\in C^4(\Omega)$ solve $F[u]=0$, where $F\in C^2(\Gamma)$ is elliptic
with respect to $u$, satisfies (17.43), and is concave in $r$ on the range of
$D^2u$.  For every $\Omega'\Subset\Omega$,
\begin{equation}\tag{17.52}
[D^2u]_{\alpha;\Omega'}\le C.
\end{equation}
Here $\alpha$ depends only on $n,\lambda,\Lambda$, and $C$ also depends on
$|u|_{2;\Omega}$, $\operatorname{dist}(\Omega',\partial\Omega)$, and the
first and second derivatives of $F$ other than $F_{rr}$.  The same conclusion
holds when $F$ is convex in $r$.
\end{theorem}
\begin{proof}
\sketch Apply the weak Harnack inequality to the nonnegative deficits of the
directional second derivatives in
Lemma~\ref{gt-ch17-lem-finite-directional-decomposition}.  Concavity supplies the reverse
control through the equation itself. For $h_k=D_{\gamma_k\gamma_k}u$,
$v=\sum_kh_k^2$, and $w_k=h_k+\varepsilon v$, the differentiated equation
and Cauchy's inequality yield
\begin{equation}
\tag{17.48}
F_{ij}D_{ij}w_k\geq-\lambda\bar\mu,
\end{equation}
with $\bar\mu$ controlled by the stated derivatives of $F$. If
$W_{ks}=\sup_{B_{sR}}w_k$, weak Harnack gives
\begin{equation}
\tag{17.49}
\Phi_{p,R}(W_{k2}-w_k)
\leq C(W_{k2}-W_{k1}+\bar\mu R^2),
\end{equation}
where $\Phi_{p,R}(v)=(|B_R|^{-1}\int_{B_R}|v|^p)^{1/p}$.
Summing over $k\ne l$ gives
\begin{equation}
\tag{17.50}
\Phi_{p,R}\left(\sum_{k\ne l}(M_{k2}-h_k)\right)
\leq C((1+\varepsilon)\omega(2R)-\omega(R)+\bar\mu R^2).
\end{equation}
Summing over all fixed directions then gives
an oscillation decay
\[
\operatorname{osc}_{B_R}D^2u
\le C\left(\frac R{R_0}\right)^\alpha(1+M_2)
(1+\widetilde\mu R_0+\widetilde\mu R_0^2),
\tag{17.51}
\]
for concentric $B_R\subset B_{R_0}\Subset\Omega$, with
$M_2=\sup|D^2u|$ and $\widetilde\mu$ controlled by the stated derivatives of
$F$.  Covering $\Omega'$ yields (17.52).  Replace $F$ by $-F$ in the convex
case.
\end{proof}

The global structure conditions used for interpolation are
\begin{equation}\tag{17.53}
\begin{aligned}
0<\lambda|\xi|^2&\le F_{ij}(x,z,p,r)\xi_i\xi_j\le\Lambda|\xi|^2,\\
|F_p|,|F_z|,|F_{rx}|,|F_{px}|,|F_{zx}|&\le\mu\lambda,\\
|F_x|,|F_{xx}|&\le\mu\lambda(1+|p|+|r|),
\end{aligned}
\end{equation}
for nonzero $\xi$ and $(x,z,p,r)\in\Gamma$, where $\lambda$ is
nonincreasing and $\Lambda,\mu$ are nondecreasing functions of $|z|$.

\begin{theorem}[Interior Hessian estimate under natural structure]
\label{gt-ch17-thm-interior-hessian-natural-structure}
\dcref{ch17:17.15}\group{chapter-17}\source{gilbarg-trudinger:page-0474}
Let $u\in C^4(\Omega)$ solve $F[u]=0$ in $\Omega\subset\mathbb R^n$.
Suppose $F$ is concave or convex jointly in $z,p,r$ and satisfies (17.53).
Then
\begin{equation}\tag{17.54}
|u|^*_{2,\alpha;\Omega}\le C,
\end{equation}
where $\alpha>0$ depends only on $n$ and $\Lambda(M)/\lambda(M)$, and $C$
also depends on $\mu(M)$, $\operatorname{diam}\Omega$, and
$M=|u|_{0;\Omega}$.
\end{theorem}
\begin{proof}
\sketch Joint concavity removes the $F_{rp},F_{rz},F_{pp},F_{pz},F_{zz}$
terms from the differentiated inequality.  Conditions (17.53) and Theorem
17.14 give $|u|^*_{2,\alpha}\le C(1+|u|^*_2)$; interior interpolation absorbs
$|u|^*_2$ into the zeroth-order bound.
\end{proof}

\section*{17.5. Dirichlet Problem for Uniformly Elliptic Equations}

\begin{lemma}[Interior regularity by difference quotients]
\label{gt-ch17-lem-interior-regularity-difference-quotients}
\dcref{ch17:17.16}\group{chapter-17}\source{gilbarg-trudinger:page-0475}
Let $u\in C^2(\Omega)$ solve $F[u]=0$, with $F$ elliptic with respect to $u$.
If $F\in C^k(\Gamma)$, $k\ge1$, then
$u\in W^{k+2,p}_{\mathrm{loc}}(\Omega)$ for every $p<\infty$.  If
$F\in C^{k,\alpha}(\Gamma)$, $0<\alpha<1$, then
$u\in C^{k+2,\alpha}(\Omega)$.
\end{lemma}
\begin{proof}
\sketch A coordinate difference quotient $w=\Delta_\ell^hu$ solves
\begin{equation}\tag{17.56}
a^{ij}D_{ij}w+b^iD_iw+cw=-f,
\end{equation}
where the coefficients are averages of $F_{ij},F_{p_i},F_z$ along the
translated segment and $f$ is the corresponding average of $F_{x_\ell}$.
Interior $L^p$ estimates, followed by Sobolev embedding, yield
$W^{3,p}_{\mathrm{loc}}$ and $C^{2,\beta}_{\mathrm{loc}}$ regularity.
Schauder estimates then give $C^{3,\alpha}$ regularity, and iteration proves
the general statement.
\end{proof}

\begin{theorem}[Dirichlet problem for concave uniformly elliptic equations]
\label{gt-ch17-thm-dirichlet-concave-uniformly-elliptic}
\dcref{ch17:17.17}\group{chapter-17}\source{gilbarg-trudinger:page-0476}
Let $\Omega\subset\mathbb R^n$ be bounded and satisfy an exterior sphere
condition at every boundary point.  Suppose $F\in C^2(\Gamma)$ is concave or
convex jointly in $z,p,r$, is nonincreasing in $z$, and satisfies (17.53).
For every $\phi\in C^0(\partial\Omega)$, the problem
\[
F[u]=0\quad\text{in }\Omega,
\qquad u=\phi\quad\text{on }\partial\Omega
\]
has a unique solution in $C^2(\Omega)\cap C^0(\overline\Omega)$.
\end{theorem}
\begin{proof}
\sketch Choose $\eta_m\in C_0^2(\Omega)$ with $\eta_m=1$ when
$\operatorname{dist}(x,\partial\Omega)\ge1/m$ and set
\begin{equation}\tag{17.57}
F_m[u]=\eta_mF[u]+(1-\eta_m)\Delta u.
\end{equation}
For smooth boundary data, solve
\begin{equation}
\tag{17.58}
F_m[u]=0\quad\text{in }\Omega,
\qquad u=\phi\quad\text{on }\partial\Omega,
\end{equation}
by applying
Theorem~\ref{gt-ch17-thm-dirichlet-continuity} to
\[
F_m[u]-tF_m[\phi]=0,
\qquad u=\phi\text{ on }\partial\Omega,
\qquad 0\le t\le1.
\tag{17.59}
\]
Theorems~\ref{gt-ch17-thm-17-3} and
\ref{gt-ch17-thm-interior-evans-krylov} give uniform zeroth-order and interior
$C^{2,\alpha}$ estimates, while $F_m=\Delta$ near the boundary gives the
global estimate.  Let $m\to\infty$ using
Theorem~\ref{gt-ch17-thm-interior-hessian-natural-structure} and boundary
barriers.
Approximate continuous data and exterior-sphere domains; uniqueness follows
from comparison.
\end{proof}

More generally, let $F^1,\ldots,F^m$ satisfy the hypotheses used in
Theorem~\ref{gt-ch17-thm-interior-evans-krylov}, and let
$G\in C^2(\mathbb R^m)$ be concave with
\begin{equation}\tag{17.60}
1\le\sum_{\nu=1}^mD_\nu G\le K.
\end{equation}
Then $G(F^1[u],\ldots,F^m[u])$ has the same interior estimate, with the
structural constants enlarged by at most the factor $K$; the corresponding
existence argument also applies.

\begin{theorem}[Dirichlet problem for countable Bellman families]
\label{gt-ch17-thm-dirichlet-countable-bellman-family}
\dcref{ch17:17.18}\group{chapter-17}\source{gilbarg-trudinger:page-0478}
Let $\Omega\subset\mathbb R^n$ be bounded and satisfy an exterior sphere
condition at every boundary point.  Suppose
$F^1,F^2,\ldots\in C^2(\Gamma)$ are concave jointly in $z,p,r$,
nonincreasing in $z$, and satisfy (17.53) uniformly.  For every
$\phi\in C^0(\partial\Omega)$, the problem
\begin{equation}\tag{17.61}
\inf_{\nu\ge1}F^\nu[u]=0\quad\text{in }\Omega,
\qquad u=\phi\quad\text{on }\partial\Omega
\end{equation}
has a unique solution in $C^2(\Omega)\cap C^0(\overline\Omega)$.
\end{theorem}
\begin{proof}
\sketch For finitely many operators, mollify
$G_0(y)=\min_\nu y_\nu$ to concave maps $G_h$ satisfying
$\sum_\nu D_\nu G_h=1$.  The estimates and existence proof for
$G_h(F^1[u],\ldots,F^m[u])$ are uniform in $h$ and $m$.  First let $h\to0$,
then exhaust the countable family and use compactness.  Comparison gives
uniqueness.
\end{proof}

In particular, Theorem~\ref{gt-ch17-thm-dirichlet-countable-bellman-family}
covers the Bellman equation (17.8) for a
countable $V$ when, uniformly in $\nu$,
\begin{equation}\tag{17.62}
\lambda|\xi|^2\le a_\nu^{ij}\xi_i\xi_j\le\Lambda|\xi|^2,
\qquad
|a_\nu^{ij}|_{2;\Omega},|b_\nu^i|_{2;\Omega},|c_\nu|_{2;\Omega},
|f_\nu|_{2;\Omega}\le\mu\lambda,
\qquad c_\nu\le0.
\end{equation}
It also covers a separable metric index set when the coefficients depend
continuously on $\nu$ pointwise in $x$, and it covers (17.6) when
$f\in C^2(\Omega)\cap L^\infty(\Omega)$.

\section*{17.6. Second Derivative Estimates for Monge--Amp\`ere Type Equations}

Consider a convex solution of
\begin{equation}\tag{17.63}
\det D^2u=f(x,u,Du)>0.
\end{equation}
Writing
\begin{equation}\tag{17.64}
\log\det D^2u=g(x,u,Du),
\qquad g=\log f,
\end{equation}
gives $F_{ij}=u^{ij}$ and
\begin{equation}
\tag{17.65}
F_{ij}=u^{ij},
\qquad F_{ij,kl}=-u^{ik}u^{jl}=-F_{ik}F_{jl},
\end{equation}
so $F$ is concave on the positive cone.

\begin{theorem}[Second derivative bounds for Monge--Amp\`ere type equations]
\label{gt-ch17-thm-monge-ampere-second-derivative-bounds}
\dcref{ch17:17.19}\group{chapter-17}\source{gilbarg-trudinger:page-0482}
Let $u\in C^2(\Omega)$ be a convex solution of (17.63), where
$f\in C^2(\Omega\times\mathbb R\times\mathbb R^n)$ is positive and
\begin{equation}\tag{17.68}
|Dg(x,u,Du)|,\ |D^2g(x,u,Du)|\le\mu.
\end{equation}
If $u\in C^2(\overline\Omega)$, then
\begin{equation}\tag{17.71}
\sup_\Omega|D^2u|\le C,
\end{equation}
where $C$ depends on $n,\mu,|u|_{1;\Omega}$ and
$\sup_{\partial\Omega}|D^2u|$.  If instead
$u\in C^{0,1}(\overline\Omega)$ is constant on $\partial\Omega$, then for
every $\Omega'\Subset\Omega$,
\begin{equation}\tag{17.72}
\sup_{\Omega'}|D^2u|\le\frac C{d_{\Omega'}},
\qquad d_{\Omega'}=\operatorname{dist}(\Omega',\partial\Omega),
\end{equation}
where $C$ depends on $n,\mu,|u|_{1;\Omega}$ and
$\operatorname{diam}\Omega$.
\end{theorem}
\begin{proof}
\sketch For a unit vector $\gamma$, apply the maximum principle to
$w=\eta e^{\beta|Du|^2/2}D_{\gamma\gamma}u$. Twice differentiating (17.64)
gives
\begin{equation}
\tag{17.66}
F_{ij}D_{ij\gamma\gamma}u
=F_{ik}F_{jl}D_{ij\gamma}uD_{kl\gamma}u+D_{\gamma\gamma}g.
\end{equation}
Consequently,
\begin{equation}
\tag{17.67}
\begin{aligned}
(\eta h)^{-1}F_{ij}D_{ij}w
&\geq D_{\gamma\gamma}u\left(
\frac{F_{ij}D_{ij}\eta}{\eta}
-\frac{F_{ij}D_i\eta D_j\eta}{\eta^2}
+(\log h)_{p_kp_l}F_{ij}D_{ik}uD_{jl}u\right.\\
&\hspace{8em}\left.+(\log h)_{p_k}F_{ij}D_{ijk}u\right).
\end{aligned}
\end{equation}
At a maximizing direction the positive quadratic term
controls the remaining third-derivative term.  Take $\eta=1$ for the global
estimate, obtaining
\begin{equation}
\tag{17.69}
\sup_\Omega w\leq C(n,\mu,|u|_{1;\Omega}).
\end{equation}
For the interior estimate, normalize the boundary constant to
zero and take $\eta=-u$; convexity gives
\[
\frac{u(x)}{\operatorname{dist}(x,\partial\Omega)}
\le\frac{\inf_\Omega u}{\operatorname{diam}\Omega},
\tag{17.70}
\]
which converts the weighted bound into (17.72).
\end{proof}

\begin{theorem}[Global Hessian bound for the zero-boundary Monge--Amp\`ere problem]
\label{gt-ch17-thm-global-hessian-zero-boundary-monge-ampere}
\dcref{ch17:17.20}\group{chapter-17}\source{gilbarg-trudinger:page-0483}
Let $\Omega$ be uniformly convex with $\partial\Omega\in C^3$, and let
$u\in C^3(\overline\Omega)$ be a convex solution of (17.64) satisfying
$u=0$ on $\partial\Omega$.  Assume
$f\in C^2(\overline\Omega\times\mathbb R\times\mathbb R^n)$ is positive.
Then
\begin{equation}\tag{17.75}
\sup_\Omega|D^2u|\le C,
\end{equation}
where $C$ depends on $n$, $|u|_{1;\Omega}$, $f$, and $\partial\Omega$.
\end{theorem}
\begin{proof}
\sketch Boundary flattening and the differentiated equation bound the mixed
tangential-normal derivatives. By (17.65), the leading coordinate-change term
in (17.26) simplifies to
\begin{equation}
\tag{17.73}
2F_{ij}D_i\left(\frac{\partial x_l}{\partial y_k}\right)D_{jl}u
=2D_i\left(\frac{\partial x_i}{\partial y_k}\right).
\end{equation}
In principal boundary coordinates, the
Monge--Amp\`ere equation reads
\[
\det D^2u=|D_nu|^{n-2}\prod_{i=1}^{n-1}\kappa_i
\left(D_nuD_{nn}u-\sum_{i=1}^{n-1}\frac{(D_{in}u)^2}{\kappa_i}\right),
\tag{17.74}
\]
so uniform convexity bounds $D_{nn}u$ on the boundary.  Apply
Theorem~\ref{gt-ch17-thm-monge-ampere-second-derivative-bounds}.
\end{proof}

\section*{17.7. Dirichlet Problem for Monge--Amp\`ere Type Equations}

For convex $u$,
\begin{equation}\tag{17.76}
\sup_\Omega|Du|=\sup_{\partial\Omega}|Du|.
\end{equation}
Assume near $\partial\Omega$ that, whenever $|p|\ge\mu(|z|)$,
\begin{equation}\tag{17.77}
0\le f(x,z,p)\le\mu(|z|)d_x^\beta|p|^\gamma,
\qquad d_x=\operatorname{dist}(x,\partial\Omega),
\qquad \beta=\gamma-n-1\ge0,
\end{equation}
with $\mu$ nondecreasing.

\begin{theorem}[Boundary gradient estimate for convex solutions]
\label{gt-ch17-thm-boundary-gradient-convex-solutions}
\dcref{ch17:17.21}\group{chapter-17}\source{gilbarg-trudinger:page-0484}
Let $\Omega$ be uniformly convex and let
$u\in C^0(\overline\Omega)\cap C^2(\Omega)$ and
$\phi\in C^2(\Omega)\cap C^0(\overline\Omega)$ be convex.  Suppose
\begin{equation}\tag{17.78}
\det D^2u=f(x,u,Du)\quad\text{in }\Omega,
\qquad u=\phi\quad\text{on }\partial\Omega,
\end{equation}
and (17.77) holds in a boundary neighbourhood $\mathcal N$.  Then
\begin{equation}\tag{17.79}
\sup_\Omega|Du|\le C,
\end{equation}
where $C$ depends on $n,\mu,\beta,\mathcal N,\Omega,|u|_{0;\Omega}$, and
$|\phi|_{1;\Omega}$.
\end{theorem}
\begin{proof}
\sketch At each boundary point use an enclosing sphere and a lower barrier
$w=\phi-\psi(d)$, choosing $\psi$ so that
$\det D^2w\ge f(x,u,Dw)$ under (17.77).  Comparison bounds the inward
difference quotients of $u$ at the boundary.  Convexity supplies the opposite
bound, and (17.76) propagates the estimate through $\Omega$.
\end{proof}

\begin{theorem}[Two-dimensional Monge--Amp\`ere Dirichlet problem]
\label{gt-ch17-thm-two-dimensional-monge-ampere-dirichlet}
\dcref{ch17:17.22}\group{chapter-17}\source{gilbarg-trudinger:page-0485}
Let $\Omega\subset\mathbb R^2$ be uniformly convex with
$\partial\Omega\in C^3$.  Suppose
$f\in C^2(\overline\Omega\times\mathbb R\times\mathbb R^2)$ is positive,
$f_z\ge0$, and (17.13), (17.14), and (17.77) hold with $n=2$ and $\beta=0$.
Then
\begin{equation}\tag{17.80}
\det D^2u=f(x,u,Du)\quad\text{in }\Omega,
\qquad u=0\quad\text{on }\partial\Omega
\end{equation}
has a solution $u\in C^{2,\alpha}(\overline\Omega)$ for every $\alpha<1$.
\end{theorem}
\begin{proof}
\sketch Apply Theorem~\ref{gt-ch17-thm-dirichlet-continuity} to
$\det D^2u/f(x,u,Du)-1=\sigma F[\phi]$ with a uniformly convex
zero-boundary function $\phi$.  Theorems~\ref{gt-ch17-thm-17-4},
\ref{gt-ch17-thm-global-hessian-holder-two-dimensional},
\ref{gt-ch17-thm-global-hessian-zero-boundary-monge-ampere}, and
\ref{gt-ch17-thm-boundary-gradient-convex-solutions} give
the required uniform bounds.  If $g_\infty<\infty$, choose $\phi$ so that
$F[\phi]\le0$ to preserve the mass bound along the homotopy.
\end{proof}

\begin{theorem}[Higher-dimensional smooth Monge--Amp\`ere Dirichlet problem]
\label{gt-ch17-thm-higher-dimensional-smooth-monge-ampere-dirichlet}
\dcref{ch17:17.23}\group{chapter-17}\source{gilbarg-trudinger:page-0486}
Let $\Omega\subset\mathbb R^n$ be uniformly convex with
$\partial\Omega\in C^4$.  Suppose
$f\in C^2(\overline\Omega\times\mathbb R\times\mathbb R^n)$ is positive,
$f_z\ge0$, and (17.13), (17.14), and (17.77) hold with $\beta=0$.
Then (17.80) has a solution $u\in C^{3,\alpha}(\overline\Omega)$ for every
$\alpha<1$.
\end{theorem}
\begin{proof}
\sketch Use the same continuity path as in
Theorem~\ref{gt-ch17-thm-two-dimensional-monge-ampere-dirichlet}.  The global
second-derivative H\"older estimate of
Theorem~\ref{gt-ch17-thm-global-third-derivative-strict-concavity} replaces the
two-dimensional estimate, and Schauder regularity gives $C^{3,\alpha}$.
\end{proof}

The Hessian bound in
Theorem~\ref{gt-ch17-thm-global-hessian-zero-boundary-monge-ampere} and the
existence conclusions in
Theorems~\ref{gt-ch17-thm-two-dimensional-monge-ampere-dirichlet} and
\ref{gt-ch17-thm-higher-dimensional-smooth-monge-ampere-dirichlet} extend to
general boundary data in $C^4(\overline\Omega)$.

\begin{theorem}[Interiorly smooth Monge--Amp\`ere Dirichlet problem]
\label{gt-ch17-thm-interior-smooth-monge-ampere-dirichlet}
\dcref{ch17:17.24}\group{chapter-17}\source{gilbarg-trudinger:page-0486}
Let $\Omega\subset\mathbb R^n$ be uniformly convex.  Suppose
$f\in C^2(\Omega\times\mathbb R\times\mathbb R^n)$ is positive,
$f_z\ge0$, and (17.13), (17.14), and (17.77) hold.  Then (17.80) has a
solution
\[
u\in C^{3,\alpha}(\Omega)\cap C^{0,1}(\overline\Omega)
\qquad\text{for every }\alpha<1.
\]
\end{theorem}
\begin{proof}
\sketch Choose bounded $f_m\le f$ agreeing with $f$ on
$|z|+|p|\le m$, and exhaust $\Omega$ by smooth uniformly convex subdomains.
Theorems~\ref{gt-ch17-thm-two-dimensional-monge-ampere-dirichlet} and
\ref{gt-ch17-thm-higher-dimensional-smooth-monge-ampere-dirichlet} solve the
approximating problems.  Theorems~\ref{gt-ch17-thm-17-4},
\ref{gt-ch17-thm-interior-evans-krylov},
\ref{gt-ch17-thm-monge-ampere-second-derivative-bounds}, and
\ref{gt-ch17-thm-boundary-gradient-convex-solutions} give compactness on
every interior subdomain and a
uniform boundary Lipschitz bound.  For large $m$, the a priori bounds keep
$(u_m,Du_m)$ in the region where $f_m=f$.
\end{proof}

\begin{corollary}[Prescribed Gauss curvature for convex graphs]
\label{gt-ch17-cor-prescribed-gauss-curvature-convex-graphs}
\dcref{ch17:17.25}\group{chapter-17}\source{gilbarg-trudinger:page-0486}
Let $\Omega\subset\mathbb R^n$ be uniformly convex and let
$K\in C^2(\Omega)\cap C^{0,1}(\overline\Omega)$ be positive in $\Omega$ and
satisfy
\begin{equation}\tag{17.81}
K=0\quad\text{on }\partial\Omega,
\qquad \int_\Omega K\,dx<\omega_n.
\end{equation}
There is a unique convex
$u\in C^2(\Omega)\cap C^{0,1}(\overline\Omega)$, vanishing on
$\partial\Omega$, whose graph has Gauss curvature $K(x)$.
\end{corollary}
\begin{proof}
\sketch Apply Theorem~\ref{gt-ch17-thm-interior-smooth-monge-ampere-dirichlet} to
$f(x,z,p)=K(x)(1+|p|^2)^{(n+2)/2}$; uniqueness follows from comparison.
\end{proof}

The strict mass condition is also necessary for a boundary-Lipschitz convex
solution.  Indeed, if
\begin{equation}\tag{17.82}
f(x,z,p)\ge\frac{h(x)}{g(p)}
\end{equation}
with positive $h\in L^1(\Omega)$ and
$g\in L^1_{\mathrm{loc}}(\mathbb R^n)$, then the injective normal map
$Du$ gives
\begin{equation}\tag{17.83}
\int_\Omega h\le g_\infty:=\int_{\mathbb R^n}g(p)\,dp.
\end{equation}
If $Du(\Omega)\ne\mathbb R^n$, in particular when
$u\in C^{0,1}(\overline\Omega)$, one must have
\begin{equation}\tag{17.84}
\int_\Omega h<g_\infty.
\end{equation}
For prescribed Gauss curvature with arbitrary nonzero boundary data in
$C^2(\overline\Omega)$, the boundary condition $K=0$ is both necessary and
sufficient for the corresponding extension of
Corollary~\ref{gt-ch17-cor-prescribed-gauss-curvature-convex-graphs}.

\section*{17.8. Global Second Derivative H\"older Estimates}

Strengthen concavity along $u$ by assuming that, for every symmetric matrix
$\xi$,
\begin{equation}\tag{17.85}
-F_{ij,kl}(x,u,Du,D^2u)\xi_{ij}\xi_{kl}\ge\lambda_0|\xi|^2
\end{equation}
for some $\lambda_0>0$.

\begin{theorem}[Global third derivative estimate under strict concavity]
\label{gt-ch17-thm-global-third-derivative-strict-concavity}
\dcref{ch17:17.26}\group{chapter-17}\source{gilbarg-trudinger:page-0488}
Let $\Omega$ be bounded with $\partial\Omega\in C^4$ and let
$\phi\in C^4(\overline\Omega)$.  Suppose
$u\in C^3(\overline\Omega)\cap C^4(\Omega)$ solves
$F[u]=0$ in $\Omega$, $u=\phi$ on $\partial\Omega$, where
$F\in C^2(\overline\Gamma)$ is elliptic with respect to $u$, satisfies
(17.43), is concave in $r$ on the range of $D^2u$, and satisfies (17.85).
Then
\begin{equation}\tag{17.86}
\sup_\Omega|D^3u|\le C,
\end{equation}
where $C$ depends on $n,\lambda,\Lambda,\lambda_0,\partial\Omega$,
$|\phi|_{4;\Omega}$, $|u|_{2;\Omega}$, and the first and second derivatives
of $F$.
\end{theorem}
\begin{proof}
\sketch Strict concavity turns the twice-differentiated equation into
$F_{ij}D_{ij}D_{\gamma\gamma}u\ge-C$; more precisely,
\begin{equation}
\tag{17.87}
F_{ij}D_{ij\gamma\gamma}u
\geq-A_{ij\gamma}D_{ij\gamma}u-B_\gamma
+\lambda_0|D^2D_\gamma u|^2\geq-C.
\end{equation}
An exterior-sphere barrier $w$ gives
\begin{equation}
\tag{17.88}
D_{\gamma\gamma}u(x)-D_{\gamma\gamma}u(x_0)
\leq\varepsilon+Cw(x)+2\sup_{\partial\Omega}|D_{\gamma\gamma}u|
\frac{|x-x_0|^2}{\delta^2}
\leq\varepsilon+C|x-x_0|.
\end{equation}
The directional decomposition and the equation imply
\begin{equation}
\tag{17.89}
|D^2u(x)-D^2u(x_0)|\leq\varepsilon+C|x-x_0|.
\end{equation}
Thus the boundary modulus of $D^2u$ reduces to that of its trace. After
flattening the boundary, tangential directions give a one-sided bound on
$D_{\gamma\gamma n}u$:
\begin{equation}
\tag{17.90}
D_{\gamma\gamma n}u(x_0)\geq-C.
\end{equation}
Apply
Lemma~\ref{gt-ch17-lem-logarithmic-boundary-modulus-convex-traces} to
$h=D_nu+k|x|^2$ to control mixed derivatives on the boundary; the equation
then controls $D_{nn}u$.  Boundary $L^p$ estimates and Morrey's inequality,
together with Theorem~\ref{gt-ch17-thm-interior-evans-krylov}, give a global
$C^{2,\alpha}$ estimate.  Schauder
estimates for (17.23) yield (17.86).
\end{proof}

\begin{lemma}[Logarithmic boundary modulus for convex traces]
\label{gt-ch17-lem-logarithmic-boundary-modulus-convex-traces}
\dcref{ch17:17.27}\group{chapter-17}\source{gilbarg-trudinger:page-0490}
Let $B^+$ be a half-ball with flat boundary portion $T$, and let
$h\in C^2(B^+)\cap C^0(B^+\cup T)$ satisfy
\begin{equation}\tag{17.91}
a^{ij}D_{ij}h\le f
\end{equation}
in $B^+$.  Assume $[a^{ij}]$ is uniformly elliptic, $f/\lambda$ is bounded,
and $h|_T$ is convex.  For $x,y\in T$ in the concentric half-radius disk and
$i=1,\ldots,n-1$,
\begin{equation}\tag{17.92}
|D_ih(x)-D_ih(y)|\le\frac{C}{1+|\log|x-y||},
\end{equation}
where $C$ depends only on $n$, $\sup\Lambda/\lambda$,
$\sup f/\lambda$, the half-ball radius, and $\sup|Dh|$.
\end{lemma}
\begin{proof}
\sketch Subtract a supporting affine function at one boundary point and
normalize
\begin{equation}
\tag{17.93}
h(0)=0=D_ih(0),\qquad i=1,\ldots,n-1.
\end{equation}
Rotate so that $D_1h(y)=a\geq0$ and the other tangential derivatives vanish.
It suffices to prove
\begin{equation}
\tag{17.94}
a\leq\frac C{|\log|y||}.
\end{equation}
Convexity gives
\begin{equation}
\tag{17.95}
h(x',0)\geq h(y',0)+a(x_1-y_1)=a(x_1-\beta),
\end{equation}
where
\begin{equation}
\tag{17.96}
0\leq\beta\leq y_1.
\end{equation}
Use the barrier
\begin{equation}
\tag{17.97}
\begin{aligned}
w(x',x_n)&=\frac a2\sqrt{(x_1-\beta)^2+x_n^2}
+\frac a2(x_1-\beta)\\
&\quad-a\gamma x_n\log\sqrt{(x_1-\beta)^2+x_n^2}
-D(x_n+|x'|^2-Ex_n^2).
\end{aligned}
\end{equation}
The constants can be chosen so that
\begin{equation}
\tag{17.98}
Lw\geq f\geq Lh\qquad\text{in }B_\varepsilon^+,
\end{equation}
and
\begin{equation}
\tag{17.99}
w\leq h\qquad\text{on }\partial B_\varepsilon^+.
\end{equation}
The maximum principle and the normal derivative at the origin give (17.94);
boundedness of $Dh$ handles non-small separations.
\end{proof}

\begin{theorem}[Global Hessian H\"older estimate without strict concavity]
\label{gt-ch17-thm-global-hessian-holder-no-strict-concavity}
\dcref{ch17:17.26'}\group{chapter-17}\source{gilbarg-trudinger:page-0493}
Let $\Omega$ be bounded with $\partial\Omega\in C^3$ and let
$\phi\in C^3(\overline\Omega)$.  Suppose
$u\in C^3(\overline\Omega)\cap C^4(\Omega)$ solves
$F[u]=0$ in $\Omega$, $u=\phi$ on $\partial\Omega$, where
$F\in C^2(\Gamma)$ is elliptic with respect to $u$, satisfies (17.43), and
is concave in $r$ on the range of $D^2u$.  Then
\begin{equation}\tag{17.86'}
|u|_{2,\alpha;\Omega}\le C,
\end{equation}
where $\alpha$ depends only on $n,\lambda,\Lambda$, and $C$ also depends on
$|u|_{2;\Omega}$ and the first and second derivatives of $F$ other than
$F_{rr}$.
\end{theorem}
\begin{proof}
\sketch After flattening the boundary, apply the boundary H\"older estimate
for uniformly elliptic equations directly to the equations satisfied by the
tangential derivatives $D_ku$.  This controls the mixed second derivatives
without a modulus for the coefficients of (17.23).  Alternatively, regard
$D_{\gamma\gamma}u$ as a function of both $x$ and the tangential direction
$\gamma$, add a uniformly elliptic operator in the $\gamma$ variables, and
extend the boundary barrier quadratically in $\gamma$.  This recovers the
one-sided third-derivative estimate without (17.85), after which the argument
of Theorem~\ref{gt-ch17-thm-global-third-derivative-strict-concavity} applies.
\end{proof}

\begin{theorem}[Global regularity for Bellman families]
\label{gt-ch17-thm-global-regularity-bellman-families}
\dcref{ch17:17.18'}\group{chapter-17}\source{gilbarg-trudinger:page-0494}
Let $\Omega$ be bounded with $\partial\Omega\in C^3$ and
$\phi\in C^3(\overline\Omega)$.  Suppose
$F^1,F^2,\ldots\in C^2(\overline\Gamma)$ are concave jointly in $z,p,r$,
nonincreasing in $z$, and satisfy (17.53) uniformly.  Then (17.61) has a
unique solution $u\in C^{2,\alpha}(\overline\Omega)$ for some $\alpha>0$
depending only on $n,\lambda,\Lambda$.
\end{theorem}
\begin{proof}
\sketch Apply the smooth finite-family approximation used for
Theorem~\ref{gt-ch17-thm-dirichlet-countable-bellman-family}.
Theorem~\ref{gt-ch17-thm-global-hessian-holder-no-strict-concavity} and
interpolation replace the dependence on
$|u|_{2;\Omega}$ by a uniform zeroth-order bound, so compactness passes the
global estimate to the countable infimum.  Comparison gives uniqueness.
\end{proof}

\section*{17.9. Nonlinear Boundary Value Problems}

Let
\begin{equation}\tag{17.100}
F[u]=F(x,u,Du,D^2u)=0\quad\text{in }\Omega,
\qquad G[u]=G(x,u,Du)=0\quad\text{on }\partial\Omega.
\end{equation}
The choice
\begin{equation}
\tag{17.101}
G(x,z,p)=z-\phi(x)
\end{equation}
is the Dirichlet condition. Otherwise, writing $\nu$ for the outer unit
normal, $G$ is oblique at $(x,z,p)$ when
\begin{equation}\tag{17.102}
G_p(x,z,p)\cdot\nu(x)>0,
\end{equation}
and oblique with respect to $u$ when
\begin{equation}
\tag{17.103}
G_p(x,u(x),Du(x))\cdot\nu(x)>0
\qquad(x\in\partial\Omega).
\end{equation}
For
$P[u]=(F[u],G[u])$, the derivative is
\begin{equation}\tag{17.104}
P_uh=(Lh,Nh),
\qquad Nh=G_zh+G_{p_i}D_i h,
\end{equation}
where $L$ is given by (17.22).

\begin{theorem}[Continuity method for nonlinear boundary operators]
\label{gt-ch17-thm-continuity-nonlinear-boundary-operators}
\dcref{ch17:17.28}\group{chapter-17}\source{gilbarg-trudinger:page-0495}
Let $\Omega$ be a $C^{2,\alpha}$ domain, $0<\alpha<1$, let
$\mathcal U\subset C^{2,\alpha}(\overline\Omega)$ be open, and fix
$\psi\in\mathcal U$.  Set
\begin{equation}\tag{17.105}
E=\{u\in\mathcal U:P[u]=\sigma P[\psi]
\text{ for some }\sigma\in[0,1]\}.
\end{equation}
Assume $F\in C^{2,\alpha}(\overline\Gamma)$,
$G\in C^{3,\alpha}(\Gamma')$, and
\begin{enumerate}
\item[(i)] $F$ is strictly elliptic with respect to every $u\in E$;
\item[(ii)] for every $u\in E$, $G$ is either Dirichlet,
$G(x,z,p)=z-\phi(x)$, or is oblique with respect to $u$;
\item[(iii)] $F_z(x,u,Du,D^2u)\le0$ and $G_z(x,u,Du)\ge0$ for every
$u\in E$, and the two functions are not both identically zero;
\item[(iv)] $E$ is bounded in $C^{2,\alpha}(\overline\Omega)$;
\item[(v)] $\overline E\subset\mathcal U$.
\end{enumerate}
Then (17.100) has a solution in $\mathcal U$.
\end{theorem}
\begin{proof}
\sketch Linear oblique Schauder theory makes (17.104) invertible under
(i)--(iii).  Theorem~\ref{gt-ch17-thm-17-6} makes the continuation set open;
(iv)--(v) and compactness make it closed.
Theorem~\ref{gt-ch17-thm-17-7} finishes the homotopy.
\end{proof}

For the remaining results let
\begin{equation}\tag{17.106}
Qu=a^{ij}(x,u,Du)D_{ij}u+b(x,u,Du).
\end{equation}

\begin{lemma}[Closedness for quasilinear oblique problems]
\label{gt-ch17-lem-closedness-quasilinear-oblique-problems}
\dcref{ch17:17.29}\group{chapter-17}\source{gilbarg-trudinger:page-0496}
Let $Q$ have coefficients $a^{ij},b\in C^1(\Omega\times\mathbb R\times
\mathbb R^n)$, and suppose $G,G_p\in C^{1,\alpha}(\partial\Omega\times
\mathbb R\times\mathbb R^n)$, $0<\alpha<1$.  Assume $Q$ is strictly
elliptic and $G$ is oblique.  Let
$u_m\in C^{2,\alpha}(\overline\Omega)$,
$f_m\in C^\alpha(\overline\Omega)$, and
$\phi_m\in C^{1,\alpha}(\partial\Omega)$ satisfy
\[
P[u_m]=(f_m,\phi_m).
\]
Suppose these three sequences are uniformly bounded, respectively, in
$C^{1,\alpha}(\overline\Omega)$, $C^\alpha(\overline\Omega)$, and
$C^{1,\alpha}(\partial\Omega)$, and converge uniformly to $u,f,\phi$.
Then $u\in C^{2,\alpha}(\overline\Omega)$ and $P[u]=(f,\phi)$.
\end{lemma}
\begin{proof}
\sketch The difference $v=u_m-u_k$ solves a linear uniformly elliptic
oblique problem.  Schauder estimates at exponent $\alpha\delta$, combined
with interpolation, give
\begin{equation}\tag{17.107}
[v]_{2,\alpha\delta}\le C+\left(\frac14+\varepsilon(m,k)\right)K_2,
\end{equation}
where $K_2=\max([D^2u_m]_{\alpha\delta},[D^2u_k]_{\alpha\delta})$ and
$\varepsilon(m,k)\to0$.  If these seminorms were unbounded, choose $m,k$ so
that
\begin{equation}
\tag{17.108}
K_2=[u_m]_{2,\alpha\delta}>4[u_k]_{2,\alpha\delta}\geq4C,
\qquad\varepsilon(m,k)<\frac14.
\end{equation}
Then (17.107) gives a
contradiction.  Compactness yields $u\in C^{2,\alpha\delta}$ and the limit
equation, and a second Schauder step raises the exponent to $\alpha$.
\end{proof}

\begin{theorem}[Quasilinear oblique solvability from $C^{1,\delta}$ bounds]
\label{gt-ch17-thm-quasilinear-oblique-solvability-c1delta}
\dcref{ch17:17.30}\group{chapter-17}\source{gilbarg-trudinger:page-0497}
Let $\Omega$ be a $C^{2,\alpha}$ domain, $0<\alpha<1$.  Suppose $Q$ is
strictly elliptic with $a^{ij},b\in C^2(\overline\Omega\times\mathbb R
\times\mathbb R^n)$, and
$G\in C^3(\partial\Omega\times\mathbb R\times\mathbb R^n)$ is oblique.
Assume
\[
a^{ij}_z=0,
\qquad b_z\le0,
\qquad G_z\ge0,
\]
and, for every $u\in C^{2,\alpha}(\overline\Omega)$, at least one of
$b_z(x,u,Du)$ or $G_z(x,u,Du)$ is nonzero somewhere on its respective
domain.  Fix $\psi\in C^{2,\alpha}(\overline\Omega)$ and set
\begin{equation}\tag{17.109}
E=\left\{u\in C^{2,\alpha}(\overline\Omega):
Qu=\sigma Q\psi\text{ in }\Omega,
\ G[u]=\sigma G[\psi]\text{ on }\partial\Omega
\text{ for some }\sigma\in[0,1]\right\}.
\end{equation}
If $E$ is bounded in $C^{1,\delta}(\overline\Omega)$ for some $\delta>0$,
then
\[
Qu=0\quad\text{in }\Omega,
\qquad G[u]=0\quad\text{on }\partial\Omega
\]
has a unique solution in $C^{2,\alpha}(\overline\Omega)$.
\end{theorem}
\begin{proof}
\sketch The implicit-function argument makes the continuation set open.
Any convergent parameter sequence has solutions uniformly controlled in
$C^{1,\delta}$; interpolation supplies the hypotheses of
Lemma~\ref{gt-ch17-lem-closedness-quasilinear-oblique-problems}, which
gives a $C^{2,\alpha}$ solution at the limiting parameter.  Thus the set is
closed and Theorem~\ref{gt-ch17-thm-17-7} gives existence.  The sign and
nondegeneracy
conditions yield comparison and uniqueness.
\end{proof}

One alternative homotopy is
\begin{equation}
\tag{17.110}
\begin{aligned}
Q_\sigma u&=a^{ij}(x,u,Du)D_{ij}u+\sigma b(x,u,Du)=0
&&\text{in }\Omega,\\
G_\sigma u&=\sigma G(x,u,Du)+(1-\sigma)u=0
&&\text{on }\partial\Omega.
\end{aligned}
\end{equation}
A typical conormal problem is
\begin{equation}
\tag{17.111}
\begin{aligned}
Qu&=\operatorname{div}A(x,u,Du)+B(x,u,Du)=0
&&\text{in }\Omega,\\
G[u]&=A(x,u,Du)\cdot\nu(x)+\phi(x,u)=0
&&\text{on }\partial\Omega.
\end{aligned}
\end{equation}
